% !TeX root = ../main.tex
\chapter{Implementacja}
\label{chap:implementation}

% FIRST DRAFT. The course of the work follows the phase write-ups in docs/
% (phase-0-*, phase-1-*), the ADRs and the performance results. Listings are
% read straight from the repository: after merging main into inzynierka,
% check that the line ranges still point at the same code.

Ten rozdział opisuje, jak zbudowaliśmy system, i nie pomija tego, co po drodze nie działało.
Wiele decyzji z poprzedniego rozdziału wygląda na oczywiste dopiero z perspektywy problemów, które je wymusiły.
Zaczynamy od struktury repozytorium i konwencji, następnie opisujemy infrastrukturę, pozyskiwanie danych, odtwarzanie strumienia, potok przetwarzania, model językowy ze środowiskiem \ac{mlops} i warstwę raportową.

\section{Struktura repozytorium i konwencje}
\label{sec:impl-repository}

Cały system mieszka w jednym publicznym repozytorium~\cite{popvibecheck}.
Kod usług napisaliśmy w języku Python 3.12, infrastrukturę w języku HCL narzędzia Terraform, a widoki raportowe w dialekcie SQL narzędzia Dataform.
Najważniejsze katalogi zestawia tabela~\ref{tab:repo-layout}.

\begin{table}[hbp]
  \caption{Najważniejsze katalogi repozytorium}
  \label{tab:repo-layout}
  \centering
  \small
  \begin{tabular}{|l|p{0.62\textwidth}|}
    \hline
    Katalog                        & Zawartość                                                                                         \\
    \hline
    \texttt{collectors/}           & wspólna biblioteka kolektorów oraz kolektory YouTube, Steam i Reddit wraz z plikami konfiguracji \\
    \hline
    \texttt{publisher/}            & wydawca odtworzeń: ładowanie archiwum do hurtowni i chronologiczna publikacja strumienia          \\
    \hline
    \texttt{dataflow/}             & potok Apache Beam, obraz szablonu Flex Template, skrypty uruchomienia, scalania i weryfikacji     \\
    \hline
    \texttt{nlp/}                  & interfejs klasyfikatora, zaślepka, klient punktu końcowego, trenowanie, ocena, integracja z MLflow \\
    \hline
    \texttt{definitions/}, \texttt{includes/} & projekt Dataform: widoki raportowe, testy danych i wspólne fragmenty SQL              \\
    \hline
    \texttt{terraform/}            & przygotowanie projektu, kompozycje środowisk \texttt{dev} i \texttt{prod} oraz 15 modułów        \\
    \hline
    \texttt{mlflow\_server/}       & obraz serwera śledzenia MLflow                                                                    \\
    \hline
    \texttt{perf/}                 & narzędzia do badań wydajności i kosztów oraz ich wyniki                                           \\
    \hline
    \texttt{docs/}                 & opisy kolejnych faz projektu i zapisy decyzji architektonicznych                                  \\
    \hline
  \end{tabular}
\end{table}

Od początku przyjęliśmy kilka zasad, których pilnują narzędzia, a nie dobra wola.
Kod formatuje \code{black}, sprawdza \code{ruff} razem z regułami dokumentowania, a typy weryfikuje \code{mypy} w trybie ścisłym.
Każda zmiana trafia do gałęzi głównej przez zgłoszenie zmian z przeglądem, z opisem w konwencji Conventional Commits~\cite{conventionalcommits}.
Decyzje, które zmieniają ustalenia projektowe, zapisujemy jako \ac{adr}; w repozytorium są cztery: zastąpienie serwisu Reddit recenzjami Steam, odtworzenie jako jedno uruchomienie w Cloud Build, przeniesienie widoków raportowych do Dataform i serwer MLflow w Cloud Run.
Każda faza projektu kończyła się dokumentem opisującym, co zbudowano, co zweryfikowano i co nie zadziałało; z tych dokumentów pochodzi większość opisanych niżej faktów.

\section{Infrastruktura i ciągłe wdrażanie}
\label{sec:impl-infrastructure}

Infrastrukturę opisuje 15 modułów Terraform, złożonych w kompozycji środowiska \code{terraform/envs/dev}.
Moduły odpowiadają rodzajom zasobów: sieci, kontom usług, magazynom, zbiorom BigQuery, tematom Pub/Sub, zadaniom Cloud Run, potokowi Dataflow, środowisku Agent Platform, serwerowi MLflow, projektowi Dataform, wyzwalaczom Cloud Build, alarmom i budżetowi.
Nazwę każdego zasobu poprzedza przedrostek, w studium przypadku \code{co}, od tytułu gry, a kończy nazwa środowiska, na przykład \code{co-raw-archive-dev}.
Dla kolejnej premiery wystarczy inny przedrostek, a Terraform utworzy obok drugi, niezależny komplet zasobów.

Stan Terraform przechowujemy w zasobniku Cloud Storage, a zmiany stosuje wyłącznie potok w Cloud Build, działający jako dedykowane konto usługi.
Laptopy członków zespołu wykonują tylko plan zmian.
Wartości zmiennych środowiska, w tym adresy osób z dostępem do raportu, przechowujemy w Secret Manager, a nie w plikach na czyimkolwiek dysku.

Potok \ac{cicd} ma dwie ścieżki.
Dla zgłoszenia zmian Cloud Build uruchamia testy jednostkowe, \code{ruff} i \code{black}, sprawdzenie plików Terraform narzędziami \code{tflint} i \code{tfsec}, budowę obrazów i plan zmian infrastruktury.
Po scaleniu buduje obrazy, oznacza je pełnym skrótem wersji kodu, publikuje w Artifact Registry i wdraża do zadań Cloud Run według skrótu obrazu, a nie etykiety, więc zadanie zawsze uruchamia dokładnie ten obraz, który powstał z danej wersji.
Zmiany infrastruktury czekają na zatwierdzenie przez człowieka w konsoli Cloud Build.
Te same narzędzia, w tych samych wersjach, uruchamiają się lokalnie jako zaczepy przed zatwierdzeniem zmian.

Do tej postaci doszliśmy etapami.
Przez pierwsze miesiące Cloud Build w ogóle nie działał: ręcznie zlecane budowanie uruchamiało się jako domyślne konto usługi Compute Engine, które w naszym projekcie nie miało żadnych ról i nie mogło odczytać nawet własnego zasobnika ze źródłami.
Obrazy kolektora, wydawcy i potoku budowaliśmy więc lokalnie w Dockerze i publikowaliśmy ręcznie, aż powstał osobny moduł z kontem usługi i wyzwalaczami.
Podobne lekcje dawało konto, którym działa Terraform.
Dodanie reguły zapory sieciowej dla potoku wymagało roli administratora bezpieczeństwa sieci, a zdefiniowanie własnej roli dla Pub/Sub, o której piszemy w podrozdziale~\ref{sec:impl-pipeline}, roli administratora ról.

Koszty pilnują alarmy i budżet.
Moduł monitorowania wysyła powiadomienia e-mail o błędzie zadania potoku, o zadaniu działającym dłużej niż 6 godzin, o komunikatach czekających w kolejce dłużej niż 30 minut i o każdym komunikacie w kolejce błędów.
Budżet miesięczny środowiska wynosił 50~zł, z powiadomieniami po przekroczeniu 50, 90, 100 i \qty{120}{\percent} tej kwoty.
Budżet okazał się mniej szczelny, niż sądziliśmy.
Przez pewien czas był przypięty do innego konta rozliczeniowego niż projekt i raportował zerowe wydatki.
Z kolei maszyna środowiska Workbench naliczała opłaty za dyski także wtedy, gdy była zatrzymana: od 10 września do 3 listopada 2026~r. kosztowało to około 140~zł miesięcznie.
Pierwszy problem naprawiliśmy w Terraform, drugi usuwając maszynę po skopiowaniu jej zawartości do Cloud Storage.
\todo{Uzupełnić łączny koszt projektu w chmurze (z rozliczeń konta).}

\section{Pozyskiwanie danych}
\label{sec:impl-collection}

Kolektory są zadaniami Cloud Run Jobs, uruchamianymi przez operatora z parametrami: identyfikatorem wydarzenia i oknem czasowym.
Wspólna biblioteka \code{collectors/common} zawiera budowę rekordu we wspólnym formacie, pseudonimizację autora, zapis paczek do archiwum, odczyt konfiguracji wydarzeń i ponawianie zapytań z wykładniczo rosnącym odstępem dla błędów przeciążenia i błędów serwera.
Konfiguracja wydarzeń jest częścią obrazu, więc uruchomienie z nieznanym identyfikatorem wydarzenia kończy się błędem, zanim kolektor wyśle choćby jedno zapytanie.

Pseudonimizację przedstawia listing~\ref{lst:author-hash}.
Funkcja odmawia działania, gdy brakuje soli, bo skrót bez soli dałoby się odwrócić, sprawdzając popularne nazwy użytkowników.

\lstinputlisting[
  language=Python, firstline=18, lastline=50, firstnumber=18, float=hbp,
  caption={Pseudonimizacja autora opinii (plik \texttt{collectors/common/author\_hash.py})},
  label={lst:author-hash},
]{../collectors/common/author_hash.py}

\subsection{Kolektor YouTube}
\label{sec:impl-youtube}

Kolektor YouTube przegląda filmy przypisane do wydarzenia, pobiera dla każdego wątki komentarzy wraz z odpowiedziami i zostawia tylko te opublikowane w oknie czasowym.
Interfejs zwraca komentarze bez względu na datę, więc filtrowanie po czasie odbywa się po stronie kolektora.
Pierwsze uruchomienie w chmurze, 5 lipca 2026~r., poprzedziły dwa problemy.
Klucz \ac{api} zapisany w Secret Manager był nieprawidłowy, miał 28 znaków zamiast 39, i serwis odrzucał go komunikatem \code{API_KEY_INVALID}; utworzyliśmy nowy klucz, ograniczony wyłącznie do YouTube Data API, a stary wyłączyliśmy.
Drugim problemem był wspomniany już brak działającego Cloud Build.

Pierwsze cztery przebiegi zebrały \num{4228} komentarzy z czterech wydarzeń.
Próbka rekordów potwierdziła zgodność ze schematem i poprawność pseudonimów, ale ujawniła też problem, który wróci w kilku miejscach pracy.
Biblioteka \code{langdetect}~\cite{langdetect}, której używamy do wykrywania języka, oznaczyła krótki angielski komentarz pisany wielkimi literami jako niemiecki.
Na krótkich tekstach wykrywanie języka jest zawodne, więc język zapisany przez kolektor traktujemy wyłącznie jako wstępny, a ostateczny ustala potok.
Po rozszerzeniu listy filmów na pełny cykl życia gry zbiór YouTube urósł do \num{12254} komentarzy z dziewięciu wydarzeń, od 9 czerwca 2024~r. do 30 kwietnia 2026~r.

\subsection{Kolektor Steam}
\label{sec:impl-steam}

Kolektor Steam jest prostszy, bo nie potrzebuje klucza ani biblioteki klienckiej.
Pobiera recenzje stronami po sto, z przerwą \qty{0.5}{\second} między zapytaniami, przedstawia się w nagłówku żądania nazwą projektu i ogranicza liczbę stron w jednym przebiegu.
Każdą recenzję przypisuje do najbliższego wydarzenia, jeśli powstała w promieniu trzech dni od niego, a pozostałe oznacza jako ogólne.
Zapisuje też rekomendację autora w polu \code{voted_up}, które dodaliśmy do schematu specjalnie dla tego źródła.

Okazało się, że jedno przejście nie wystarczy, i to z powodu, który leży po stronie serwisu, a nie kolektora.
Steam kończy stronicowanie strumienia wyników długo przed jego końcem.
Przejście przez wszystkie recenzje we wszystkich językach zatrzymało się samo na stronie \num{1243}, po \num{124187} recenzjach, docierając tylko do września 2025~r., podczas gdy gra ukazała się w kwietniu.
Serwis zwrócił po prostu pustą stronę; nic nie zgłosiło błędu, a podniesienie limitu stron niczego nie zmieniało.
Węższy strumień, na przykład recenzje tylko w jednym języku, mieści się w tym limicie i sięga do premiery.
Zbiór zebraliśmy więc w czternastu przejściach, zestawionych w tabeli~\ref{tab:steam-passes}, a nakładanie się przejść nie stanowi problemu, bo scalanie w hurtowni zostawia jeden wiersz na identyfikator recenzji.

\begin{table}[hbp]
  \caption{Przejścia, z których złożono zbiór recenzji Steam}
  \label{tab:steam-passes}
  \centering
  \small
  \begin{tabular}{|p{0.62\textwidth}|r|}
    \hline
    Przejście                                                                              & Recenzji dostępnych \\
    \hline
    recenzje negatywne, wszystkie języki                                                   & \num{13409}         \\
    \hline
    angielskie, zakup w sklepie Steam                                                      & \num{103218}        \\
    \hline
    angielskie, zakup poza sklepem Steam                                                   & \num{37411}         \\
    \hline
    chiński uproszczony                                                                    & \num{39575}         \\
    \hline
    francuski                                                                              & \num{23607}         \\
    \hline
    rosyjski                                                                               & \num{17160}         \\
    \hline
    hiszpański, niemiecki, koreański, polski, chiński tradycyjny, turecki, japoński, portugalski & od \num{522} do \num{9971} każde \\
    \hline
  \end{tabular}
\end{table}

Recenzji angielskich jest \num{140654}, czyli więcej, niż mieści się w jednym strumieniu, dlatego podzieliliśmy je według sposobu zakupu, ostatniego kryterium, jakie oferuje serwis.
Dwie inne drogi sprawdziliśmy i odrzuciliśmy: parametr zakresu dni jest ograniczony do około roku, a filtr ostatnio edytowanych recenzji zwraca zaledwie 400 rekordów.
Skrypt \code{collect_history.sh} z listą przejść i uzasadnieniem każdego z nich jest w praktyce definicją zbioru danych: skoro serwis nie wydaje historii w jednym kawałku, zebrany zbiór to dokładnie suma tych strumieni.

Ostatecznie zebraliśmy \num{251639} z \num{280699} recenzji, czyli \qty{89.6}{\percent}.
Brakuje części angielskich recenzji z ograniczonych końcówek obu angielskich strumieni.
Nic nie wskazuje, by brakujące recenzje różniły się wydźwiękiem, bo recenzje negatywne zebraliśmy w całości w osobnym przejściu, ale nie jest to też próba losowa i tak ją opisujemy.

Przed odtworzeniem sprawdziliśmy jakość danych.
Nie było zduplikowanych identyfikatorów, niepoprawnych pseudonimów ani dat z przyszłości.
Wykryliśmy natomiast trzy grupy rekordów, które mogły zatrzymać przetwarzanie: 816 recenzji składających się z samych białych znaków, \num{14350} recenzji dłuższych niż \num{1000} znaków, czyli \qty{5.4}{\percent}, oraz \num{11277} recenzji o długości od jednego do trzech znaków, na przykład \enquote{gg} albo pojedyncze emoji.
Pierwsze dwie grupy wymagały zmian w kliencie modelu, opisanych w podrozdziale~\ref{sec:impl-nlp}; trzecią przetwarzamy normalnie, bo niska pewność modelu na takim tekście jest informacją, a nie szumem.

\section{Odtwarzanie strumienia}
\label{sec:impl-replay}

Wydawca odtworzeń jest zadaniem Cloud Run Job z dwoma trybami.
W trybie ładowania wczytuje wszystkie pliki archiwum do tabeli \code{raw_landing}, zastępując jej zawartość, i scala je z tabelą \code{raw_staging}, zostawiając dla każdego identyfikatora najświeżej pobraną wersję.
Ponowne ładowanie niczego nie zmienia, co sprawdziliśmy przy pierwszym uruchomieniu: po dwóch ładowaniach tabela pośrednia miała dokładnie \num{4228} wierszy.

W trybie odtworzenia wydawca czyta tabelę pośrednią posortowaną według czasu publikacji i identyfikatora i publikuje każdy rekord w temacie Pub/Sub pod kluczem porządkującym \code{global}.
Przed publikacją odczekuje czas wyliczony przez funkcję z listingu~\ref{lst:sleep}: rzeczywisty odstęp od poprzedniej opinii podzielony przez współczynnik przyspieszenia i obcięty do zadanego maksimum.
Przy domyślnym współczynniku \num{86400} jeden dzień historii trwa jedną sekundę, a obcięcie sprawia, że czas całego odtworzenia jest ograniczony przez liczbę rekordów, a nie przez długość okresu, który obejmują.

\lstinputlisting[
  language=Python, firstline=64, lastline=90, firstnumber=64, float=hbp,
  caption={Wyliczanie przerwy przed publikacją kolejnej opinii (plik \texttt{publisher/replay.py})},
  label={lst:sleep},
]{../publisher/replay.py}

Komunikat zawiera rekord dokładnie w takiej postaci, w jakiej zapisał go kolektor, a dane o odtworzeniu niesie w atrybutach: źródło, wydarzenie, oryginalny czas publikacji i czas wysłania.
Zadanie ma wyłączone ponawianie, bo ponowione częściowo wykonane odtworzenie opublikowałoby część rekordów dwukrotnie.
Klient Pub/Sub łączy się z regionalnym punktem dostępowym usługi, czego wymaga porządkowanie komunikatów.

Pierwsza wersja obrazu wydawcy nie uruchomiła się wcale.
Import dowolnego modułu wspólnej biblioteki kolektorów wykonywał kod inicjalizujący pakiet, który od razu importował moduł zapisu do Cloud Storage, a biblioteki Cloud Storage w obrazie wydawcy nie było.
Rozwiązaniem był leniwy import tego modułu.
Pierwsze odtworzenie, 16 lipca 2026~r., opublikowało \num{4228} rekordów w \qty{96.4}{\second}, a sprawdzenie wszystkich komunikatów w subskrypcji kontrolnej nie wykazało ani jednego naruszenia kolejności.

\section{Potok przetwarzania strumieniowego}
\label{sec:impl-pipeline}

Potok podzieliliśmy na dwie warstwy.
Moduł \code{dataflow/transforms.py} zawiera całą logikę na poziomie rekordu jako zwykłe funkcje, bez zależności od Apache Beam: odczyt komunikatu, walidację, wykrywanie języka i budowę wiersza wynikowego.
Moduł \code{dataflow/pipeline.py} tylko składa z nich graf przetwarzania, przedstawiony na listingu~\ref{lst:pipeline}.
Dzięki temu logikę testujemy bez uruchamiania potoku, a warstwę Beam da się przeczytać w kilka minut.

\lstinputlisting[
  language=Python, firstline=230, lastline=262, firstnumber=230, float=hbp,
  caption={Graf przetwarzania potoku (plik \texttt{dataflow/pipeline.py})},
  label={lst:pipeline},
]{../dataflow/pipeline.py}

Potok odczytuje komunikaty z subskrypcji, biorąc czas zdarzenia z atrybutu \code{created_utc}.
Rekordy, których nie da się odczytać albo którym brakuje identyfikatora lub czasu publikacji, trafiają do tematu błędnych komunikatów wraz z opisem przyczyny.
Rekord, który model uzna za trudny, nie jest błędny i trafia do hurtowni jak każdy inny, bo gubienie opinii z powodu niepewności modelu po cichu zniekształcałoby analizę.
Poprawne rekordy są grupowane w paczki po najwyżej 200 i klasyfikowane razem, a klasyfikator powstaje raz na proces roboczy, w metodzie \code{setup}, a nie dla każdego rekordu.
Wyniki trafiają do tabeli \code{events_landing} trybem dopisywania, a wiersze odrzucone przez BigQuery również kierujemy do tematu błędnych komunikatów.

\subsection{Cztery nieudane uruchomienia}
\label{sec:impl-launch-failures}

Potok nie uruchomił się za pierwszym, drugim ani trzecim razem.
Każda porażka wskazała rzeczywistą lukę, a każda poprawka jest dziś pokryta testem albo zapisana w Terraform.

Pierwsze uruchomienie przerwał błąd zapisu do BigQuery przez interfejs Storage Write API, który wymaga jawnie podanego schematu tabeli i nie odczytuje go sam.
Zamiast powielać schemat w kodzie Python, gdzie mógłby rozjechać się z tabelą zarządzaną przez Terraform, potok odczytuje go z tabeli docelowej w chwili budowania grafu.
Błąd wyszedł późno, bo testy jednostkowe sprawdzały pojedyncze transformacje, a nic nie składało całego grafu; od tego czasu mamy testy budowy grafu.

Drugie uruchomienie zakończyło się komunikatem, że do użycia tej transformacji potrzebna jest Java.
W bibliotece Beam dla języka Python zapis przez Storage Write API jest transformacją międzyjęzykową, która uruchamia usługę w Javie przy budowie grafu i wymaga środowiska Java na maszynach roboczych.
W obrazie, który z założenia ma być samowystarczalny, bo maszyny nie mają dostępu do internetu, nie chcieliśmy tego środowiska.
Domyślną metodą zapisu został więc zapis strumieniowy, w całości napisany w Pythonie.
Oszczędność, jaką daje Storage Write API, przy naszej skali nie ma znaczenia, bo całe odtworzenie to kilka megabajtów danych.

Trzecie uruchomienie było najbardziej podstępne.
Zadanie osiągnęło stan działania, raportowało, że wszystko jest w porządku, i nie przetwarzało niczego.
Zaległość w subskrypcji stała na \num{4228} komunikatach przez ponad godzinę.
Jedyną wskazówką było ostrzeżenie, które brzmiało jak uwaga o wydajności: odczyt konfiguracji subskrypcji nie powiódł się z powodu braku uprawnień.
Rola subskrybenta Pub/Sub pozwala pobierać komunikaty, ale nie odczytywać konfiguracji subskrypcji, a Dataflow robi to przed rozpoczęciem pracy.
Dodaliśmy rolę przeglądającego.

Czwarte uruchomienie ujawniło prawdziwą blokadę.
Ponieważ potok bierze czas zdarzenia z atrybutu komunikatu, Dataflow tworzy na temacie źródłowym własną, pomocniczą subskrypcję do śledzenia znacznika wodnego i usuwa ją po zatrzymaniu zadania.
Wymaga to uprawnień do tworzenia subskrypcji i dołączania ich do tematu.
Dokumentacja zaleca rolę edytora Pub/Sub dla całego projektu, która pozwala też publikować i usuwać dowolne tematy.
Zdefiniowaliśmy zamiast niej własną rolę z dokładnie pięcioma potrzebnymi uprawnieniami.
Oba ostatnie problemy łączy to, że zawodzą po cichu: zadanie działa, skaluje się, raportuje dobry stan i nie robi nic.
Od tamtej pory pierwszą rzeczą, którą sprawdzamy, jest zaległość w subskrypcji, a nie stan zadania.

\subsection{Błąd powtarzalności w wykrywaniu języka}
\label{sec:impl-langdetect}

Gdy potok wreszcie działał, dwa pierwsze odtworzenia tych samych danych dały różne wyniki.
Liczba wierszy, identyfikatory, wydźwięk i wszystkie inne kolumny były identyczne, ale 68 rekordów różniło się językiem.
Ponieważ język jest w tabeli zdarzeń wartością ostateczną, już samo to łamało obietnicę powtarzalności, dla której zbudowaliśmy całą architekturę.

Przyczyny były dwie i obie tkwiły w bibliotece \code{langdetect}.
Biblioteka losuje fragmenty tekstu i dla każdego wykrycia tworzy nowy generator liczb losowych z ziarnem zapisanym w atrybucie klasy, domyślnie pustym, co oznacza ziarno z entropii systemu.
Ustawialiśmy ziarno przy imporcie modułu, co wyglądało na równoważne ustawianiu go przy każdym wywołaniu, a nie było: na maszynach roboczych moduł docierał do detektora bez tego przypisania.
Krótki tekst \enquote{Hell yeah!} bez ziarna bywał rozpoznawany jako angielski, indonezyjski albo turecki.
Ustawienie ziarna w samej funkcji wykrywania, widoczne na listingu~\ref{lst:langdetect}, usunęło część różnic, ale nie wszystkie: po poprawce różniło się wciąż 87 rekordów.

\lstinputlisting[
  language=Python, firstline=232, lastline=257, firstnumber=232, float=hbp,
  caption={Deterministyczne wykrywanie języka (plik \texttt{dataflow/transforms.py})},
  label={lst:langdetect},
]{../dataflow/transforms.py}

Wzorzec różnic zdradził drugą przyczynę.
Na świeżo uruchomionych maszynach zwykły angielski tekst dostawał oznaczenia \code{af}, \code{it} albo \code{no}, a na rozgrzanych \code{en}.
Biblioteka przy pierwszym użyciu tworzy globalny obiekt detektora i dopiero potem wczytuje do niego profile języków.
Wątek, który trafi w tę lukę, widzi obiekt, uznaje go za gotowy i wykrywa język na podstawie profili wczytanych do tej pory, czyli alfabetycznie pierwszych, stąd afrikaans.
Dataflow uruchamia kilka wątków przetwarzania na każdej maszynie, więc na zimnej maszynie dzieje się to niemal przy każdym starcie.
Odtworzyliśmy to lokalnie: zimny start i dwanaście równoległych wątków dawały błędne języki w trzech próbach na trzy.
Pierwsze wczytanie profili zabezpieczyliśmy blokadą i przenieśliśmy do metody \code{setup}, zanim jakikolwiek wątek dotrze do detektora, a testy odtwarzają zimny stan biblioteki i obciążają ją dwunastoma wątkami.

Z tego błędu wynika ogólniejsza lekcja.
Błąd nie tkwił w naszym kodzie, tylko w zachowaniu cudzej biblioteki przy współbieżności, którą narzuca Dataflow.
Każdy komponent wczytujący pliki przy starcie zasługuje na taką samą podejrzliwość, a odcisk tabeli jest narzędziem, które takie problemy wychwytuje.

\subsection{Scalanie i odcisk tabeli}
\label{sec:impl-merge}

Scalanie wyników z tabelą zdarzeń wykonuje instrukcja z listingu~\ref{lst:merge}.
Z tabeli lądowania wybiera dla każdego identyfikatora wiersz z najpóźniejszym czasem przetworzenia i albo aktualizuje nim istniejący wiersz, jeśli jest świeższy, albo go wstawia.
Czasu pobrania nie można użyć do rozstrzygania, bo dwa odtworzenia tych samych danych mają go identyczny.

\lstinputlisting[
  language=SQL, firstline=14, lastline=43, firstnumber=14, float=hbp,
  caption={Scalanie tabeli lądowania z tabelą zdarzeń (plik \texttt{dataflow/promote.sql})},
  label={lst:merge},
]{../dataflow/promote.sql}

Odcisk tabeli, przedstawiony na listingu~\ref{lst:fingerprint}, sumuje skróty wszystkich wierszy z pominięciem wersji modelu i czasu przetworzenia.
Suma nie zależy od kolejności wierszy.
Przy pierwszym uruchomieniu zapytanie się nie powiodło, bo funkcja skrótu zwraca liczby z pełnego zakresu typu 64-bitowego i ich suma dla kilku tysięcy wierszy przekracza ten zakres; rzutowanie na typ \code{NUMERIC} rozwiązało problem.

\lstinputlisting[
  language=SQL, firstline=19, lastline=32, firstnumber=19, float=hbp,
  caption={Odcisk tabeli zdarzeń (plik \texttt{dataflow/verify.sql})},
  label={lst:fingerprint},
]{../dataflow/verify.sql}

\subsection{Odtworzenie jako jedno uruchomienie}
\label{sec:impl-replay-build}

Na początku odtworzenie było ciągiem ręcznie wpisywanych poleceń: uruchom szablon, uruchom wydawcę, obserwuj zadanie, zatrzymaj je, scal wyniki.
Dwa z tych kroków kosztowały nas prawdziwe przebiegi.
Zakończenie pracy wydawcy oznacza, że wszystko zostało opublikowane, a nie przetworzone, więc scalanie w tym momencie obcina koniec odtworzenia.
Gorzej, że przedwczesne zatrzymanie potoku porzuca wszystko, co zostało w kolejce.

Dokładnie to zdarzyło się przy pierwszym odtworzeniu z modelem XLM-R.
Pętla oczekująca uznawała, że przetwarzanie się skończyło, gdy liczba wierszy w tabeli lądowania nie zmieniła się przez trzy minuty, a zablokowany potok wygląda dokładnie tak samo.
Przetworzono \num{10002} z \num{12254} rekordów, w subskrypcji zostało \num{3190} komunikatów, a przebieg zgłosił sukces.
Tabela zdarzeń wyszła mieszana: \num{10002} wiersze z etykietami nowego modelu i \num{2252} z etykietami zaślepki z wcześniejszego odtworzenia.
Sprawdzenie kompletności tego nie wychwyciło, bo porównywało identyfikatory z tabeli pośredniej ze wszystkimi kiedykolwiek wylądowanymi, a te identyfikatory wylądowały wcześniej.

Odpowiedzią była decyzja, że jedno odtworzenie to jedno uruchomienie wyzwalacza w Cloud Build, z regułami opisanymi w podrozdziale~\ref{sec:design-replay}: czekanie na kompletność bieżącego przebiegu, brak scalania przy niekompletności i zatrzymanie potoku zawsze.
Przy wdrażaniu tej zmiany wyszły trzy kolejne błędy.
Dwa wynikały z tego samego nieporozumienia: zmienne Cloud Build nie są zmiennymi powłoki, więc zapis z podwójnym znakiem dolara rozwijał się do pustego napisu, co dawało niepoprawną nazwę zadania i zapytanie bez projektu.
Trzeci był groźniejszy: świeżo uruchomione zadanie przez kilka minut nie pojawia się na liście zadań, więc krok zatrzymujący, który szukał go po nazwie, zgłosiłby brak czegokolwiek do zatrzymania, a zadanie naliczałoby opłaty dalej.
Orkiestrator zapisuje więc identyfikator zadania z wyniku uruchomienia i posługuje się nim we wszystkich dalszych krokach.

\section{Model językowy i środowisko MLOps}
\label{sec:impl-nlp}

Część dotycząca modelu przeszła najwięcej zmian, bo pierwotny plan zderzył się z danymi i z ograniczeniami sprzętowymi.
Jej punktem wyjścia był kontrakt z listingu~\ref{lst:classifier}: potok zna wyłącznie interfejs klasyfikatora i klasę wyniku, więc zamiana modelu jest zmianą w rejestrze, a nie w potoku.
Interfejs wymaga dwóch rzeczy: determinizmu, czyli tej samej odpowiedzi dla tego samego tekstu przy każdym wywołaniu, oraz odpowiedzi dla każdego wejścia, także pustego.

\lstinputlisting[
  language=Python, firstline=136, lastline=179, firstnumber=136, float=hbp,
  caption={Interfejs klasyfikatora wydźwięku (plik \texttt{nlp/base.py})},
  label={lst:classifier},
]{../nlp/base.py}

Rejestr udostępnia dwie implementacje: deterministyczną zaślepkę, na której zbudowaliśmy i przetestowaliśmy cały potok, zanim istniał jakikolwiek model, oraz klienta punktu końcowego Agent Platform.
Nieznana nazwa modelu w parametrze uruchomienia zatrzymuje potok przy starcie maszyny, zamiast po cichu oznaczyć całe odtworzenie niewłaściwym modelem.

\subsection{Pierwszy model: angielski DistilBERT}
\label{sec:impl-distilbert}

Pierwszym modelem był angielski DistilBERT w wersji bez rozróżniania wielkości liter, dostrojony do trzech klas.
Dane treningowe, około \num{222000} przykładów, pochodziły ze zbiorów publicznych: SST-2, TweetEval, Sentiment140 i GoEmotions, uzupełnionych o 220 oznaczonych komentarzy z naszego zbioru~\cite{socher2013recursive,barbieri2020tweeteval,go2009twitter,demszky2020goemotions}.
Model trenowaliśmy na maszynie \code{e2-standard-4} w środowisku Workbench, bez akceleratora graficznego.
Trzy epoki, czyli \num{75090} kroków przy paczce ośmiu przykładów, zajęły około 44,5 godziny.

Zanim to się udało, maszyna zatrzymała się w trakcie pierwszego treningu.
Nasz moduł Terraform ustawiał wtedy wyłączanie maszyny po trzech godzinach bezczynności notatnika, a proces uczenia uruchomiony przez SSH nie liczy się jako aktywność.
Usunięcie ustawienia z metadanych maszyny nie pomogło, bo usługa wyłączająca odczytała je przy starcie systemu i odliczała dalej.
Skuteczne okazało się dopiero zablokowanie samej usługi systemowej, a w Terraform wartość zero, która miała wyłączać tę funkcję, interfejs odrzucał, więc moduł po prostu pomija ten parametr.

Na publicznym zbiorze kontrolnym model osiągnął dokładność \num{0.802} i F1 makro \num{0.786}.
Na naszym zbiorze wzorcowym z 300 komentarzy YouTube wyniki były znacznie gorsze: dokładność \num{0.533} i F1 makro \num{0.472}, a dla komentarzy angielskich dokładność około \num{0.635}.
Dokładność \num{0.533} jest praktycznie równa wynikowi reguły, która każdej opinii przypisuje klasę pozytywną, bo taką klasę ma 161 z 300 komentarzy zbioru wzorcowego.
Z analizy błędów wynikły trzy przyczyny.
Pierwszą był język: angielski model dostawał komentarze, z których wiele nie było po angielsku, a recenzje Steam, które mieliśmy dopiero zebrać, były po angielsku tylko w niespełna połowie.
Drugą była różnica dziedzin: recenzje filmów i wpisy z Twittera słabo przypominają komentarze o grach, a 220 komentarzy z naszej dziedziny ginęło wśród ponad \num{220000} przykładów publicznych.
Trzecią była nierównowaga klas: w zbiorze wzorcowym było tylko 20 komentarzy negatywnych, a model często mylił neutralne z pozytywnymi.

\subsection{Zmiana na wielojęzyczny XLM-RoBERTa}
\label{sec:impl-xlmr}

Odpowiedzią na pierwszą przyczynę był model wielojęzyczny XLM-R w wersji bazowej, z około 278 milionami parametrów, wstępnie wytrenowany na stu językach~\cite{conneau2020unsupervised}.
Na dwie pozostałe odpowiedzieliśmy zmianą danych i sposobu uczenia.
Nowy zestaw danych, zestawiony w tabeli~\ref{tab:xlmr-corpus}, liczył od \num{150000} do \num{180000} przykładów i był zdominowany przez wielojęzyczne opinie w językach, które przeważają w recenzjach Steam.
Zbiory SST-2 i Sentiment140, które najbardziej odpowiadały za różnicę dziedzin, wyłączyliśmy.
Komentarze z części treningowej naszego zbioru wzorcowego powieliliśmy dziesięciokrotnie, a funkcja straty ważyła klasy odwrotnie do ich liczebności.
Najlepszy punkt kontrolny wybieraliśmy według F1 makro, a nie dokładności, co przy nierównych klasach ma znaczenie.

\begin{table}[hbp]
  \caption{Skład danych treningowych modelu XLM-R v2}
  \label{tab:xlmr-corpus}
  \centering
  \small
  \begin{tabular}{|p{0.27\textwidth}|p{0.2\textwidth}|p{0.1\textwidth}|p{0.28\textwidth}|}
    \hline
    Zbiór                                         & Przykładów           & Udział            & Uwagi                                              \\
    \hline
    clapAI MultiLingualSentiment~\cite{clapaimultilingualsentiment} & po \num{18000} dla 5 języków i po \num{6000} dla 3 & ok. \qty{62}{\percent} & en, zh, ru, fr, ko oraz de, es, ja; teksty do 512 znaków \\
    \hline
    TweetEval, wydźwięk~\cite{barbieri2020tweeteval} & \num{25000}       & ok. \qty{15}{\percent} & angielskie wpisy z mediów społecznościowych   \\
    \hline
    XLM-T, wydźwięk wielojęzyczny~\cite{barbieri2022xlmt} & ok. \num{24000} & ok. \qty{14}{\percent} & osiem języków                            \\
    \hline
    GoEmotions~\cite{demszky2020goemotions}       & \num{15000}          & ok. \qty{9}{\percent}  & komentarze z serwisu Reddit                  \\
    \hline
    zbiór wzorcowy, część treningowa              & 220 × 10             & ok. \qty{1}{\percent}  & część kontrolna nigdy nie trafia do uczenia  \\
    \hline
  \end{tabular}
\end{table}

\subsection{Trenowanie wyłącznie na procesorach}
\label{sec:impl-cpu}

Wszystkie modele trenowaliśmy i serwowaliśmy wyłącznie na procesorach (\ac{cpu}), bez procesorów graficznych (\ac{gpu}).
Na koncie rozliczeniowym projektu nie korzystaliśmy z akceleratorów, a cały projekt musiał zmieścić się w budżecie rzędu kilkudziesięciu złotych miesięcznie.
\todo{Potwierdzić i opisać dokładną przyczynę braku GPU (limit akceleratorów na koncie, koszt, zasady konta rozliczeniowego).}
Ograniczenie to ukształtowało cały proces.
Model XLM-R jest kilkukrotnie większy od DistilBERT, więc do jego trenowania użyliśmy na czas jednego przebiegu maszyny \code{n2-standard-16} z 64~GB pamięci, przy paczce ośmiu przykładów i szacowanym czasie od 25 do 45 godzin.
Po pierwszym pobraniu modelu i danych z serwisu Hugging Face kopiowaliśmy pamięć podręczną do Cloud Storage, aby kolejne sesje działały bez dostępu do internetu.
Przebieg trwający półtorej doby oznaczał też, że każdy błąd w konfiguracji kosztował dzień pracy, a maszynę trzeba było po treningu zatrzymać ręcznie, bo automatyczne wyłączanie zostało wyłączone.
\todo{Uzupełnić faktyczny czas trenowania XLM-R v2 i v2.1 (z MLflow).}

Brak akceleratora odczuliśmy jeszcze mocniej przy serwowaniu.
Jedna replika \code{n1-standard-8} klasyfikowała recenzję w około 170~ms, co dla \num{12254} komentarzy YouTube oznaczało około 35 minut, a dla całego zbioru z recenzjami Steam około 12 godzin.
Zwiększenie paczek niewiele dawało, bo czas zajmowało samo wnioskowanie, a nie narzut wywołań.
Jedynym rozwiązaniem była równoległość, a tu czekała kolejna niespodzianka.
Wdrożenie repliki dla pełnego odtworzenia kończyło się błędem przekroczenia limitu procesorów do serwowania modeli, na każdym rozmiarze maszyny, nawet z dwoma rdzeniami.
Limit dla rodziny maszyn N1 był od 10 października zajęty w całości, choć nic nie było wdrożone, a zwolnić mógł go tylko wniosek do Google.
Okazało się jednak, że Agent Platform liczy ten limit osobno dla każdej rodziny maszyn, a limity dla rodzin N2 i N2D były wolne.
Wdrożyliśmy więc ten sam model jednocześnie na maszynach \code{n2-standard-8} i \code{n2d-standard-8}, z podziałem ruchu po połowie, co dało 16 rdzeni zamiast ośmiu.
Pełne odtworzenie \num{263893} opinii trwało w ten sposób 2 godziny i 12 minut.

\subsection{Dostrajanie kontynuowane: wersja 2.1}
\label{sec:impl-xlmr-21}

Wersja v2 wyraźnie poprawiła wyniki, ale zbiór wzorcowy z 300 komentarzy, z zaledwie 80 w części kontrolnej, nie pozwalał rzetelnie porównywać kolejnych modeli.
Rozszerzyliśmy go do 950 komentarzy: 670 w części uczącej, 80 w części strojenia i 200 w części kontrolnej, w której znalazło się wszystkich 80 komentarzy z pierwotnej części kontrolnej.
Kandydatów do oznaczenia wstępnie oznaczył asystent programistyczny oparty na dużym modelu językowym, pracujący na plikach według naszych wytycznych oznaczania.
Człowiek przeglądał i zatwierdzał wszystkie wstępnie negatywne komentarze, wszystkie o niskiej pewności i wszystkie przydzielone do części strojenia i kontrolnej.
Konsekwencje tej procedury dla wiarygodności oceny omawiamy w rozdziale~\ref{chap:evaluation}.

Wersję v2.1 uzyskaliśmy, kontynuując dostrajanie wersji v2 przez trzy epoki z niską stałą uczenia, \num{1e-5}, na danych z naszej dziedziny uzupełnionych o \num{5000} przykładów publicznych, które chroniły model przed zapomnieniem.
Przed treningiem spisaliśmy warunki, które nowa wersja musi spełnić, aby zastąpić poprzednią: wyższe F1 makro na części kontrolnej, wyższe F1 dla klasy negatywnej i brak pogorszenia na komentarzach angielskich.
Wersja v2.1 spełniła wszystkie trzy, choć poprawa dla klasy negatywnej była minimalna; szczegółowe wyniki przedstawiamy w podrozdziale~\ref{sec:eval-quality}.

\subsection{Serwowanie modelu}
\label{sec:impl-serving}

Wybrany model rejestrujemy w rejestrze modeli Agent Platform i wdrażamy na punkt końcowy wyłącznie na czas odtworzenia.
Używamy gotowego kontenera Hugging Face do wnioskowania na procesorach; kontener TorchServe, który wypróbowaliśmy wcześniej, nie akceptował katalogu modelu w formacie Hugging Face.
Po odtworzeniu model jest wycofywany z punktu końcowego, a punkt końcowy usuwany, więc między odtworzeniami serwowanie nic nie kosztuje.
Całe przejście modelu od trenowania do potoku przedstawia rysunek~\ref{fig:model-lifecycle}.

\begin{figure}[hbp]
  \centering
  \begin{tikzpicture}[
      font=\footnotesize,
      box/.style={draw, rounded corners=2pt, align=center, minimum height=9mm, text width=27mm, fill=gray!8},
      store/.style={box, fill=blue!8},
      arr/.style={-{Latex[length=2mm]}, thick},
      lab/.style={font=\scriptsize, fill=white, inner sep=1pt}]
    \node[box] (wb) at (0,0) {Agent Platform\\Workbench\\(trenowanie, CPU)};
    \node[box] (ml) at (4.2,1.2) {serwer MLflow\\(Cloud Run, Cloud SQL)};
    \node[store] (gcs) at (4.2,-1.2) {wagi modeli\\(Cloud Storage)};
    \node[box] (reg) at (8.4,0) {rejestr modeli\\Agent Platform};
    \node[box] (ep) at (8.4,-2.6) {punkt końcowy\\Agent Platform};
    \node[box] (df) at (4.2,-3.8) {potok Apache Beam\\(Dataflow)};
    \draw[arr] (wb) -- node[lab, above left=0pt] {parametry, metryki} (ml);
    \draw[arr] (wb) -- node[lab, below left=0pt] {wagi} (gcs);
    \draw[arr] (gcs) -- node[lab, below right=0pt] {rejestracja} (reg);
    \draw[arr] (reg) -- node[lab, right=1pt] {wdrożenie na czas odtworzenia} (ep);
    \draw[arr, {Latex[length=2mm]}-{Latex[length=2mm]}] (df) -- node[lab, below right=0pt] {REST} (ep);
  \end{tikzpicture}
  \caption{Cykl życia modelu: od trenowania do klasyfikacji w potoku}
  \label{fig:model-lifecycle}
\end{figure}

Dwa problemy z serwowaniem ujawniły dopiero recenzje Steam.
Pierwszym była długość: XLM-R przyjmuje najwyżej 512 tokenów, a kontener serwujący nie przycina dłuższych tekstów, tylko zwraca błąd o niezgodnym rozmiarze tensora.
Takiego błędu nie naprawi żadne ponowienie, a potok ponawia paczkę w nieskończoność, więc jedna długa recenzja wystarczała, by zatrzymać całe odtworzenie.
Klient nie ma tokenizatora, bo właśnie po to model działa zdalnie, więc długość szacuje funkcja z listingu~\ref{lst:tokens}: jeden token na znak chiński, japoński lub koreański, jeden na dwa znaki cyrylicy lub greki i jeden na cztery pozostałe znaki.
Budżet 240 tokenów wyznaczyliśmy doświadczalnie, bo przy 480 punkt końcowy wciąż odrzucał część zapytań, a gdy szacunek okaże się zbyt optymistyczny, klient skraca tekst i ponawia zapytanie.

\lstinputlisting[
  language=Python, firstline=182, lastline=198, firstnumber=182, float=hbp,
  caption={Szacowanie liczby tokenów bez tokenizatora (plik \texttt{nlp/endpoint/classifier.py})},
  label={lst:tokens},
]{../nlp/endpoint/classifier.py}

Drugim problemem były recenzje złożone z samych białych znaków.
Wysyłanie ich do modelu groziło tym samym rodzajem błędu, więc klient odpowiada na nie sam: etykietą neutralną z pewnością równą dokładnie zero.
Funkcja softmax nigdy nie zwraca zera, więc taka wartość jednoznacznie oznacza, że modelu nie pytano.
Przy błędzie, którego nie da się naprawić ponowieniem, klient zgłasza wyjątek, a nie zwraca udawanej etykiety neutralnej, bo taka etykieta po cichu zatrułaby tabelę zdarzeń.

Ostatnia zmiana dotyczyła wersji modelu zapisywanej w każdym wierszu.
Początkowo był to identyfikator wdrożenia na punkcie końcowym, który zmienia się przy każdym wdrożeniu, nawet tych samych wag.
Przy jednoczesnym wdrożeniu na dwóch rodzinach maszyn jeden model miałby więc dwie wersje.
Wiersz zawiera teraz wersję z rejestru modeli, \code{xlmr-sent-v2.1}, czyli to, co faktycznie decyduje o etykietach.

\subsection{Serwer MLflow}
\label{sec:impl-mlflow}

MLflow towarzyszył trenowaniu od początku, ale przez pierwsze miesiące nie był wdrożony.
Parametry i metryki przebiegów trafiały do katalogu na dysku maszyny Workbench, a tylko artefakty do Cloud Storage.
Miało to trzy wady: historię przebiegów widział tylko ten, kto uruchomił maszynę i interfejs MLflow, historia mogła zniknąć razem z dyskiem, a sam MLflow był umową na jednej maszynie, a nie infrastrukturą jak reszta systemu.

Dlatego 28 października 2026~r. uruchomiliśmy serwer śledzenia MLflow w Cloud Run, z bazą PostgreSQL w Cloud SQL na metadane i rejestr modeli, opisany w Terraform.
Artefakty pozostały w Cloud Storage, a serwer ich nie pośredniczy, bo wagi ważą gigabajty, a Cloud Run ogranicza rozmiar żądania do 32 MiB.
Dostęp do interfejsu daje \ac{iap} osobom z tej samej listy, która uprawnia do odczytu widoków raportowych.
Dotychczasowe przebiegi przenieśliśmy na serwer skryptem, który zachowuje stary identyfikator przebiegu jako znacznik.

Ta decyzja ma swoją cenę: baza działa stale i kosztuje od 40 do 55~zł miesięcznie, czyli mniej więcej tyle, ile wynosił cały wcześniejszy budżet środowiska.
Rozważaliśmy bazę wyłączaną między użyciami i bazę SQLite na wolumenie Cloud Storage, ale pierwsza oznaczałaby platformę zwykle niedostępną, a druga ryzyko uszkodzenia danych, bo Cloud Storage nie obsługuje blokowania plików.
Uznaliśmy, że platforma, która przez większość czasu jest wyłączona, nie jest wdrożoną platformą.
Po przeniesieniu przebiegów usunęliśmy maszynę Workbench, a dane potrzebne do ponownego trenowania zostały w Cloud Storage.

\section{Warstwa raportowa i panel}
\label{sec:impl-reporting}

Pierwsze widoki raportowe zdefiniowaliśmy w kodzie Terraform, a część miar, na przykład zestawienia tygodniowe i porównanie wydarzenia z tłem, liczył kod w Pythonie w tymczasowym panelu.
Takie rozwiązanie miało trzy wady.
Zmiana klauzuli \code{GROUP BY} wymagała planu, zatwierdzenia i zastosowania zmian infrastruktury, tak samo jak zmiana konta usługi.
Nic nie testowało danych, więc widok mógł liczyć coś dwukrotnie i i tak się wdrożyć.
Miar liczonych w Pythonie nie mógł z kolei użyć Data Studio.

Wszystkie widoki przenieśliśmy więc do projektu Dataform w katalogu głównym repozytorium.
Każdy widok jest osobnym plikiem z opisem kolumn i testami danych: unikalnością klucza, warunkami na wartości i zgodnością sum z tabelą zdarzeń.
Dla zgłoszenia zmian Cloud Build kompiluje projekt, a po scaleniu zleca Dataform wykonanie wszystkich akcji i czeka na wynik; nieudany test przerywa wdrożenie.
Zmiana miary nie wymaga już zatwierdzania infrastruktury, bo przeglądem jest zgłoszenie zmian.
Nowe widoki dały na danych produkcyjnych wyniki zgodne z wcześniejszymi obliczeniami w Pythonie z tolerancją względną $10^{-9}$.

Najciekawszym widokiem jest \code{v_event_lift}, przedstawiony na listingu~\ref{lst:event-lift}, który porównuje okno wydarzenia z tłem z tego samego źródła.
Dla każdej pary wydarzenia i źródła dzieli dzienne liczniki na te z okna wydarzenia i wszystkie pozostałe, a następnie liczy różnicę wskaźnika netto i stosunek liczby opinii na dzień.
Tłem dla recenzji Steam są głównie recenzje spoza okien wydarzeń, a dla komentarzy YouTube, które zbieramy tylko wokół wydarzeń, pozostałe wydarzenia.

\lstinputlisting[
  language=SQL, firstline=37, lastline=80, firstnumber=37, float=hbp,
  caption={Efekt wydarzenia względem tła z tego samego źródła (plik \texttt{v\_event\_lift.sqlx})},
  label={lst:event-lift},
]{../definitions/reporting/v_event_lift.sqlx}

Panel jest raportem Data Studio o nazwie \enquote{Pop vibe check}, złożonym z trzech stron opisanych w podrozdziale~\ref{sec:design-ui}.
Każdy z sześciu używanych widoków ma w raporcie własne źródło danych, a zapytania wykonują się z uprawnieniami osoby oglądającej.
Wskaźnik netto i udziały klas są polami obliczanymi z sum liczników, więc są poprawne dla dowolnego wyboru źródeł i okresów.
Raport przechowuje wyniki zapytań przez 12 godzin, co przy danych zmieniających się tylko po odtworzeniu w zupełności wystarcza.
Słabością tego rozwiązania jest to, że układ raportu, pola obliczane i źródła danych istnieją tylko w Data Studio, na koncie jednego z nas, a nie w repozytorium.
W repozytorium są wyłącznie widoki, dlatego wszystkie miary celowo liczymy w nich, a raport tylko je wyświetla.
Jak z panelu i z całego systemu korzystać w praktyce, opisuje kolejny rozdział.
